\section{The frozen survivor and excess assignment}

We use the area-preserving H\'enon map, in the polynomial family introduced
by H\'enon \citep{Henon1976},
\begin{equation}\label{eq:henon}
 H_6(q,p)=(1-6q^2-p,q),
 \qquad \det DH_6=1.
\end{equation}
The certified survivor is conjugate to the mixing subshift \(\SigmaA\) with
state labels
\[
 0={--},\qquad 1={-+},\qquad 2={+-},\qquad 3={++}
\]
and adjacency matrix
\begin{equation}\label{eq:adjacency}
 A=
 \begin{pmatrix}
 1&0&1&0\\
 1&0&0&0\\
 0&1&0&1\\
 0&1&0&0
 \end{pmatrix}.
\end{equation}
The state at time \(j\) records the signs of \((q_j,q_{j-1})\).
The symbolic object, rectangles and hyperbolicity data are inherited without
refitting from HCS-P31 \citep{WangP31}.

For a primitive orbit \(\gamma\), let \(M_\gamma\in\mathrm{SL}_2\) be the
chronological monodromy and let \(\Lambda_\gamma>1\) be the modulus of its
physical unstable multiplier.  If \(f_\gamma\) is the monic minimal
polynomial of the signed unstable multiplier, define
\begin{equation}\label{eq:height}
 \mathcal H_\gamma=\log M(f_\gamma),
 \qquad
 \ell_\gamma=\log\Lambda_\gamma,
 \qquad
 \Egal_\gamma=\mathcal H_\gamma-\ell_\gamma.
\end{equation}
Reciprocity of \(f_\gamma\) splits its roots into reciprocal pairs.  The
physical pair contributes \(\ell_\gamma\); every other pair contributes a
nonnegative logarithm.  This gives \eqref{eq:split-intro} exactly
\citep{WangP48,WangP54}.

We shall use five primitive symbolic cycles:
\begin{equation}\label{eq:witness-cycles}
\begin{aligned}
 \gamma_1&=(0),
 &\gamma_3&=(0,2,1),\\
 \gamma_{4a}&=(0,0,2,1),
 &\gamma_{4b}&=(0,2,3,1),\\
 \gamma_5&=(0,0,2,3,1).
\end{aligned}
\end{equation}
Direct enumeration of \eqref{eq:adjacency} gives respectively one, zero, one,
two and two primitive cycles in periods one through five, agreeing with the
independently certified H6 orbit catalogue.  Thus \eqref{eq:witness-cycles}
does not select a best orbit from a larger same-period sample.
